\documentclass[../../main.tex]{subfiles}% 注意这里的文件路径不能用 ./main.tex ,否则用latexmk编译子文件会报错
\graphicspath{{\subfix{./image/}}{\subfix{../../image/}}} % 指定图片引用路径,后续可以直接使用图片文件名
% 如果图片存放在上级目录,则使用 ../../image/ 作为备用路径

% 例如:
% \begin{figure}[H]
% \centering
% \includegraphics[width=0.3\textwidth]{图.png}
% \caption{}
% \label{figure:图}
% \end{figure}
% 注意:上述\label{}一定要放在\caption{}之后,否则引用图片序号会只会显示??.

\begin{document}

\section{$\mathbb{R}$上的积分换元公式}

\begin{lemma}\label{lemma:实变函数--引理5.19}
设 $f(x)$ 是 $[a,b]\subseteq \mathbb{R}$ 上的实值函数, $E\subset[a,b]$. 如果$f$在$[a,b]$上Lipschitz连续,记Lipschitz常数为$M$, 则
\[
m^*(f(E)) \leqslant Mm^*(E). 
\]
\end{lemma}
\begin{proof}
将点集 $E$ 作如下分解: 对任给 $\varepsilon>0$, 作
\[
E_n = \left\{x\in E: \text{当 } [a,b] \text{ 中的点 } y \text{ 满足 } |y-x|<\frac{1}{n} \text{ 时}, \text{ 有 } |f(y)-f(x)|\leqslant(M+\varepsilon)|x-y|\right\}.
\]
显然, 有 $E_n\subset E_{n+1}$,进而$f(E_n)\subset f(E_{n+1})$.于是由\refcor{corollary:递增集列外测度和极限可交换}可知
\[
\lim_{n\to\infty}m^*(E_n) = m^*(E), \quad \lim_{n\to\infty}m^*(f(E_n)) = m^*(f(E)).
\]
从而根据 $\varepsilon$ 的任意性, 我们只需证明对每个 $n$ 有
\[
m^*(f(E_n)) < (M+\varepsilon)(m^*(E_n)+\varepsilon).
\]
为此, 对任意固定的$n\in \mathbb{N}$,在 $(a,b)$ 中取覆盖 $E_n$ 的开区间列 $\{I_{n,k}\}$, 使得
\[
\sum_{k=1}^{\infty} |I_{n,k}| < m^*(E_n)+\varepsilon, \quad |I_{n,k}| < \frac{1}{n}, \quad k=1,2,\cdots.
\]
由条件知,对$\forall k\in \mathbb{N},\forall s,t\in E_n\cap I_{n,k}$, 有
\[
|f(s)-f(t)| < (M+\varepsilon)|I_{n,k}|.
\]
于是再结合\rrefpro{proposition:外测度相关的基本性质}{proposition:外测度相关的基本性质-6}知对$\forall k\in \mathbb{N}$,有
\begin{align*}
m^*(f(E_n\cap I_{n,k}))\leqslant \operatorname{diam}(f(E_n\cap I_{n,k})) 
\leqslant (M+\varepsilon)|I_{n,k}|.
\end{align*}
故
\begin{align*}
m^*(f(E_n)) &= m^*\left(f\left(E_n\cap\bigcup_{k=1}^{\infty} I_{n,k}\right)\right) \leqslant \sum_{k=1}^{\infty} m^*(f(E_n\cap I_{n,k})) \\
&\leqslant \sum_{k=1}^{\infty} \operatorname{diam}(f(E_n\cap I_{n,k})) 
\leqslant (M+\varepsilon)\sum_{k=1}^{\infty} |I_{n,k}| < (M+\varepsilon)(m(E_n)+\varepsilon).
\end{align*}
\end{proof}

\begin{corollary}\label{corollary:可微的可测函数像集外测度小于导数积分}
若 $f(x)$ 是 $[a,b]$ 上的可测函数, $E\subset[a,b]$ 是可测集, 且 $f(x)$ 在 $E$ 上可微, 则
\[
m^*(f(E)) \leqslant \int_E |f'(x)|\mathrm{d}x. \label{5.17}
\]
\end{corollary}

\begin{proof}
易知 $f'(x)$ 是 $E$ 上的可测函数. 对任给的 $\varepsilon>0$, 作集合列
\[
E_n = \{x\in E: (n-1)\varepsilon\leqslant|f'(x)|<n\varepsilon\}, \quad n=1,2,\cdots.
\]
由\reflem{lemma:实变函数--引理5.19}, 我们有
\begin{align*}
m^*(f(E_n)) \leqslant n\varepsilon m(E_n) \leqslant (n-1)\varepsilon m(E_n) + \varepsilon m(E_n) 
\leqslant \int_{E_n} |f'(x)|\mathrm{d}x + \varepsilon m(E_n).
\end{align*}
由此可知
\begin{align*}
m^*(f(E)) \leqslant \sum_{n=1}^{\infty} m^*(f(E_n)) 
\leqslant \sum_{n=1}^{\infty} \int_{E_n} |f'(x)|\mathrm{d}x + \varepsilon \sum_{n=1}^{\infty} m(E_n) = \int_E |f'(x)|\mathrm{d}x + \varepsilon m(E).
\end{align*}
由 $\varepsilon$ 的任意性即得 $m^*(f(E))\leqslant\int_E |f'(x)|\mathrm{d}x$.
\end{proof}

\begin{proposition}
设 $f(x)$ 在 $[a,b]$ 上是处处可微的, 且 $f'(x)$ 是 $[a,b]$ 上的可积函数, 则
\[
\int_a^b f'(x)\mathrm{d}x = f(b)-f(a).
\]
\end{proposition}
\begin{proof}
因为 $f'\in L([a,b])$, 所以对于任意的 $\varepsilon>0$, 存在 $\delta>0$, 当 $e\subset[a,b]$ 且 $m(e)<\delta$ 时, 有
\[
\int_e |f'(x)|\mathrm{d}x < \varepsilon.
\]
从而对于其长度总和小于 $\delta$ 的任意的互不相交区间组 $(x_1,y_1), (x_2,y_2), \cdots, (x_n,y_n)$, 可知
\begin{align*}
\sum_{i=1}^n |f(y_i)-f(x_i)| \leqslant \sum_{i=1}^n m(f([x_i,y_i])) \leqslant \sum_{i=1}^n \int_{[x_i,y_i]} |f'(x)|\mathrm{d}x = \int_{\bigcup_{i=1}^n [x_i,y_i]} |f'(x)|\mathrm{d}x < \varepsilon.
\end{align*}
这说明 $f(x)$ 是 $[a,b]$ 上的绝对连续函数, 故结论成立.
\end{proof}

\begin{theorem}\label{theorem:实变函数--定理5.21}
设 $f(x)$ 是 $[a,b]$ 上的实值函数, 在 $[a,b]$ 的子集 $E$ 上是可微的, 则有
\begin{align*}
f'(x)=0, \text{a.e.} x\in E\iff m(f(E))=0.
\end{align*}
\end{theorem}
\begin{proof}

$\Longrightarrow$: 令 $E_n = \{x\in E: n-1\leqslant|f'(x)|<n\}$ ($n\in\mathbb{N}$), 则由\reflem{lemma:实变函数--引理5.19}知
\[
m^*(f(E)) \leqslant \sum_{n=1}^{\infty} m^*(f(E_n)) \leqslant \sum_{n=1}^{\infty} n m^*(E_n).
\]
所以由 $m^*(E_n)=0$ 可知 $m(f(E))=0$.

$\Longleftarrow$: 作点集
\[
B_n \equiv \left\{x\in E: \text{当 } |y-x|<\frac{1}{n} \text{ 时}, |f(y)-f(x)|\geqslant\frac{|y-x|}{n}\right\}.
\]
显然有
\[
B \equiv \{x\in E: |f'(x)|>0\} = \bigcup_{n=1}^{\infty} B_n.
\]
从而只需证明 $m(B_n)=0$ (一切 $n$) 即可. 为此, 又只需指出对于任一个长度小于 $\frac{1}{n}$ 的区间 $I$, 有 $m(I\cap B_n)=0$ 即可.

记 $A=I\cap B_n$, 因为由假设知 $m(f(A))=0$, 所以对任意的 $\varepsilon>0$, 存在区间列 $\{I_k\}$, 使得
\[
\bigcup_{k=1}^{\infty} I_k \supset f(A), \quad \sum_{k=1}^{\infty} |I_k| < \varepsilon.
\]
现在令 $A_k = A\cap f^{-1}(I_k)$. 因为
\[
A = \bigcup_{k=1}^{\infty} A_k, \quad \text{且} \quad A_k \subset I\cap B_n,
\]
又注意到 $f(A_k)\subset I_k$, 所以我们有
\begin{align*}
m^*(A) \leqslant \sum_{k=1}^{\infty} m^*(A_k) \leqslant \sum_{k=1}^{\infty} \operatorname{diam}(A_k) 
\leqslant \sum_{k=1}^{\infty} n \cdot \operatorname{diam}(f(A_k)) \leqslant n \sum_{k=1}^{\infty} m(I_k) \leqslant n\varepsilon.
\end{align*}
由 $\varepsilon$ 的任意性知 $m(A)=0$.
\end{proof}

\begin{theorem}[复合函数的微分]\label{theorem:实变函数--复合函数的微分}
设 $g: [a,b]\to[c,d]$ 是几乎处处可微的函数, $F(x)$ 是 $[c,d]$ 上的几乎处处可微的函数且 $F'(x)=f(x)$ a.e., $F(g(t))$ 在 $[a,b]$ 上是几乎处处可微的. 若对于 $[c,d]$ 中的任一零测集 $Z$, 总有 $m(F(Z))=0$, 则
\begin{align}
(F(g(t)))' = f(g(t))g'(t), \quad \text{a.e. } t\in[a,b]. \label{eq::-2-490238402309raesf5.18}
\end{align}
\end{theorem}
\begin{remark}
注意, 在上述定理中, $F$ 将零测集映为零测集这一条件是重要的. 事实上, 设 $g$ 是 $[0,1]$ 上的严格递增的连续函数且 $g'(t)=0$ a.e., 而令 $F = g^{-1}$, 易知 $F(x)$ 是单调且几乎处处可微的函数, $(F(g(t)))' = 1$ ($t\in[a,b]$), 但是 \eqref{eq::-2-490238402309raesf5.18}式不成立.
\end{remark}
\begin{proof}
令 $Z = \{x\in[c,d]: F \text{ 在 } x \text{ 处不可微}\}$, $A = g^{-1}(Z)$ 且 $B = [a,b]\setminus A$. 对于 $B$ 中使 $g$ 可微的点 $t$, 根据 $g$ 在 $t$ 处的连续性以及 $F'(g(t))=f(g(t))$, 我们有
\[
F(g(t+h)) - F(g(t)) = (f(g(t))+\delta(h))(g(t+h)-g(t)),
\]
其中当 $h\to 0$ 时, $\delta(h)\to 0$ (若 $g(t+h)=g(t)$, 令 $\delta(h)=0$). 用 $h$ 除上式两端, 并令 $h\to 0$, 可知\eqref{eq::-2-490238402309raesf5.18}式在 $B$ 上几乎处处成立.

因为 $m(g(A))\leqslant m(Z)=0$, 所以由假定得 $m(F(g(A)))=0$. 从而根据\refthe{theorem:实变函数--定理5.21}, 可知
\[
g'(t) = 0 = (F(g(t)))', \quad \text{a.e. } t\in A.
\]
综合上述结果, 得到
\[
(F(g(t)))' = f(g(t))g'(t), \quad \text{a.e. } t\in[a,b].
\]
\end{proof}

\begin{corollary}
设 $g(t)$ 以及 $f(g(t))$ 在 $[a,b]$ 上几乎处处可微, 其中 $f(x)$ 在 $[c,d]$ 上绝对连续, $g([a,b])\subset[c,d]$, 则
\[
(f(g(t)))' = f'(g(t))g'(t), \quad \text{a.e. } t\in[a,b].
\]
\end{corollary}

\begin{theorem}[换元积分法]\label{theorem:实变函数--换元积分法}
假设 $g(x)$ 在 $[a,b]$ 上是几乎处处可微的, $f(x)$ 是 $[c,d]$ 上的可积函数, 且 $g([a,b])\subset[c,d]$, 记
\[
F(x) = \int_c^x f(t)\mathrm{d}t,
\]
则下述两个命题是等价的:
\begin{enumerate}[(1)]
\item $F(g(t))$ 是 $[a,b]$ 上的绝对连续函数;
\item $f(g(t))g'(t)$ 是 $[a,b]$ 上的可积函数, 且有
\begin{align}
\int_{g(\alpha)}^{g(\beta)} f(x)\mathrm{d}x = \int_{\alpha}^{\beta} f(g(t))g'(t)\mathrm{d}t, \label{eq::02ir09weufjioswf5.19}
\end{align}
其中 $\alpha,\beta\in[a,b]$.
\end{enumerate}
\end{theorem}
\begin{remark}
在这个定理中, 并没有要求 $g(t)$ 是绝对连续函数. 实际上, 这也是不必要的. 例如, 令
\[
F(x) = x^2, \quad g(t) =
\begin{cases}
t\sin\frac{\pi}{x}, & 0<x\leqslant 1, \\
0, & x=0,
\end{cases}
\]
则 $g(t)$ 除在 $t=0$ 处外皆可微, 而且不是 $[0,1]$ 上的绝对连续函数. 然而, $F(x)$ 与 $F(g(t))$ 都在 $[0,1]$ 上绝对连续 (导数有界).
\end{remark}
\begin{proof}
假定 (2) 成立, 则由 \eqref{eq::02ir09weufjioswf5.19} 式知, 对一切 $\alpha,\beta\in[a,b]$, 有
\[
F(g(\beta)) - F(g(\alpha)) = \int_{g(\alpha)}^{g(\beta)} f(x)\mathrm{d}x = \int_{\alpha}^{\beta} f(g(t))g'(t)\mathrm{d}t.
\]
这说明 $F(g(t))$ 是 $[a,b]$ 上的绝对连续函数.

反之, 假定 (1) 成立, 则\hyperref[theorem:实变函数--复合函数的微分]{复合函数的微分}的条件满足. 因为 $(F(g(t)))'$ 是 $[a,b]$ 上的可积函数, 所以 $f(g(t))g'(t)$ 也是 $[a,b]$ 上的可积函数, 从而得到
\begin{align*}
\int_{g(\alpha)}^{g(\beta)} f(x)\mathrm{d}x = F(g(\beta)) - F(g(\alpha)) 
= \int_{\alpha}^{\beta} [F(g(t))]'\mathrm{d}t = \int_{\alpha}^{\beta} f(g(t))g'(t)\mathrm{d}t.
\end{align*}
\end{proof}

\begin{corollary}
设 $g: [a,b]\to[c,d]$ 是绝对连续函数, $f\in L[c,d]$, 则下述条件之一都是\eqref{eq::02ir09weufjioswf5.19}式成立的充分条件:
\begin{enumerate}[(1)]
\item $g(t)$ 在 $[a,b]$ 上是单调函数;
\item $f(x)$ 在 $[c,d]$ 上是有界函数;
\item $f(g(t))g'(t)$ 在 $[a,b]$ 上是可积函数.
\end{enumerate}
\end{corollary}
\begin{proof}
令
\[
F(x) = \int_c^x f(u)\mathrm{d}u, \quad x\in[c,d],
\]
则问题归结为证明 $F(g(t))$ 在 $[a,b]$ 上绝对连续. 对于 (3), 利用 (2) 以及\hyperref[theorem:控制收敛定理]{控制收敛定理}即可.
\end{proof}

\begin{example}
设 $f(x)$ 是 $[0,+\infty)$ 上的非负递减函数, 且对任意的 $A>0$, $f(x)$ 在 $[0,A]$ 上绝对连续, 则对 $p\geqslant 1$, 有
\[
p\int_0^{+\infty} (f(x))^p x^{p-1}\mathrm{d}x \leqslant \left(\int_0^{+\infty} f(x)\mathrm{d}x\right)^p.
\]
\end{example}
\begin{proof}
若 $f\notin L([0,\infty))$, 则上式显然成立. 若 $f\in L([0,\infty))$, 注意到 $f(x)$ 是递减的, 则 $f(x)\to 0$ ($x\to+\infty$). 从而可得
\begin{align*}
p\int_0^{+\infty} (xf(x))^{p-1}f(x)\mathrm{d}x \leqslant p\int_0^{+\infty} \left(\int_0^x f(t)\mathrm{d}t\right)^{p-1}f(x)\mathrm{d}x 
= \int_0^{+\infty} \frac{\mathrm{d}}{\mathrm{d}x}\left(\int_0^x f(t)\mathrm{d}t\right)^p \mathrm{d}x = \left(\int_0^{+\infty} f(x)\mathrm{d}x\right)^p.
\end{align*}
\end{proof}






\end{document}